\documentclass[11pt,a4paper]{article}
\usepackage[utf8]{inputenc}
\usepackage[english]{babel}
\usepackage{amsmath}
\usepackage{amsfonts}
\usepackage{booktabs}
\usepackage{hyperref}
\usepackage{listings}
\usepackage{xcolor}
\usepackage{geometry}

\geometry{margin=1in}

\definecolor{codegreen}{rgb}{0,0.6,0}
\definecolor{codegray}{rgb}{0.5,0.5,0.5}
\definecolor{backcolour}{rgb}{0.95,0.95,0.92}

\lstset{
    backgroundcolor=\color{backcolour},
    commentstyle=\color{codegreen},
    keywordstyle=\color{magenta},
    numberstyle=\tiny\color{codegray},
    stringstyle=\color{purple},
    basicstyle=\ttfamily\small,
    breakatwhitespace=false,
    breaklines=true,
    captionpos=b,
    keepspaces=true,
    numbers=left,
    numbersep=5pt,
    showspaces=false,
    showstringspaces=false,
    showtabs=false,
    tabsize=2,
    frame=single
}

\title{Advanced Macro Structure Report: Weather Maps \\ \large Dynamic Coverage Control and Shape Variation for \texttt{Clouds.hs}}
\author{Advanced Graphics Engine Group}
\date{\today}

\begin{document}

\maketitle

\section{The Core Problem: Uniform Noise Distribution}
Even with dynamic looping and multiple scattering, sampling a 3D noise texture directly using world coordinates results in an evenly distributed, repetitive grid of clouds. 3D noise functions are structurally homogeneous; they have the same average density everywhere.

To break this pattern, you must decouple the \textbf{structural detail} (the 3D noise) from the \textbf{macro distribution} (where clouds are allowed to form). This is achieved using a low-frequency \textbf{2D Weather Map}.

---

\section{Architectural Solution: The 2D Weather Map Texture}
Introduce a 2D texture wrapped over a massive scale (e.g., $20,000 \times 20,000$ meters) sampled via world horizontal positions $(x, z)$. This texture contains macroeconomic parameters packed into its color channels:

\begin{itemize}
    \item \textbf{R (Red Channel) --- Cloud Coverage:} Controls local cloud density ($0.0 = $ clear sky, $1.0 = $ heavy overcast).
    \item \textbf{G (Green Channel) --- Cloud Type / Height Mask:} Controls whether clouds are low, flat Stratus ($0.0$) or towering Cumulonimbus ($1.0$).
    \item \textbf{B (Blue Channel) --- Wetness / Storm Potential:} Used to darken cloud bellies locally during heavy rain or storms.
\end{itemize}

---

\section{Mathematical Density Modulation Shifting}
Instead of evaluating raw noise directly against a fixed threshold, pass a global uniform variable, \texttt{globalCoverage} ($0.0 \to 1.0$), from your CPU app code into the shader. Combine it with the local weather data:

\begin{equation}
    \text{CombinedCoverage} = \text{clamp}(\text{WeatherMap.R} \cdot \text{globalCoverage}, 0.0, 1.0)
\end{equation}

Modify your base density accumulation logic inside the loop using a non-linear remap operation:

\begin{equation}
    \text{FinalDensity} = \text{max}\left(0.0, \frac{\text{BaseNoise} - (1.0 - \text{CombinedCoverage})}{\text{CombinedCoverage}}\right)
\end{equation}

When \texttt{globalCoverage} is set to $0.0$, \text{CombinedCoverage} collapses to $0.0$, making the subtraction term $(1.0 - 0.0) = 1.0$, completely clipping away all noise to create a perfectly clear sky. When shifted to $1.0$, the entire noise spectrum is pushed above the threshold, forcing thick, stormy overcast clouds.

---

\section{Haskell FIR Integration Script}
Below is the updated sample block written for Haskell \texttt{FIR} syntax, illustrating how to integrate the 2D weather sampler into your structural raymarch loop.

\begin{lstlisting}[language=Haskell, caption={Haskell FIR Weather Map Masking Architecture}]
-- Update your pipeline configuration descriptors to take a 2D weather texture map
type AdvancedCloudPipeline = '[
    "weather_map"    ':-> Texture2D (Vec 4 Float),
    "cloud_noise"    ':-> Texture3D (Vec 4 Float),
    "globalCoverage" ':-> Uniform Float
  ]

evaluateWeatherDensity :: Vec3 Float -> Float -> Program AdvancedCloudPipeline Float
evaluateWeatherDensity currPos heightNorm = do
    globalCov <- use @"globalCoverage"
    
    -- 1. Sample Weather Map using massive macro landscape coordinates
    let weatherScale = 0.00005 -- Covers a large 20km area before looping
    let weatherUV = Vec2 (currPos.&X * weatherScale) (currPos.&Z * weatherScale)
    weatherSample <- sample @"weather_map" weatherUV
    
    let localCoverage = weatherSample.&R
    let cloudType     = weatherSample.&G
    
    -- 2. Dynamically modify allowed cloud thickness profile based on type
    let allowedHeightMin = lerp 0.1 0.0 cloudType -- Stratus sits lower
    let allowedHeightMax = lerp 0.4 1.0 cloudType -- Cumulus towers higher
    let heightMask = smoothstep allowedHeightMin (allowedHeightMin + 0.05) heightNorm 
                   * (1.0 - smoothstep (allowedHeightMax - 0.1) allowedHeightMax heightNorm)
    
    -- 3. Combine with base 3D noise field
    baseNoise <- sample @"cloud_noise" (currPos * 0.0002)
    
    let targetCoverage = localCoverage * globalCov
    let activeDensity  = (baseNoise.&R) - (1.0 - targetCoverage)
    
    return (max 0.0 (activeDensity * heightMask))
\end{lstlisting}

---

\section{Configuration Tuning Matrix: Clear Sky to Storms}
By updating the uniform values on the CPU frame by frame, you can easily control weather states dynamically:

\begin{table}[h!]
\centering
\small
\begin{tabular}{llll}
\toprule
\textbf{Target Weather State} & \textbf{\texttt{globalCoverage} Value} & \textbf{Visual Effect Outcome} & \textbf{Ambient Lighting Adjust} \\ \midrule
\textbf{Clear Blue Sky} & $0.00 \to 0.15$ & No cloud density survives threshold clipping. & Bright blue skylight tint. \\
\textbf{Fair Weather Cumulus} & $0.25 \to 0.45$ & Scattered, localized fluffy cloud clusters. & Normal ambient coefficients. \\
\textbf{Overcast Day} & $0.60 \to 0.80$ & Continuous cloud sheets covering the dome. & Muted grey ambient tone. \\
\textbf{Heavy Thunderstorm} & $0.85 \to 1.00$ & Tall, dark, sunlight-blocking cloud layers. & Dark blue-grey undertones. \\ \bottomrule
\end{tabular}
\caption{Global state values configuration parameter matrix.}
\end{table}

\end{document}

